\section*{Bab \ref{ch:b3:bioinorganic} --- Kimia Bioanorganik}

\begin{solution}{exo:b3:bioinorganic:1}
$\log(\theta/(1 - \theta))$ berubah dari $\log 0.25 = -0.602$ menjadi $\log 4 = +0.602$ sementara $\log p$ naik sebesar
$\log(5.9/2.1) = 0.449$: $n = 1.204/0.449 = 2.7$.
\end{solution}

\begin{solution}{exo:b3:bioinorganic:2}
$\ce{Mn^2+} < \ce{Fe^2+} < \ce{Co^2+} < \ce{Ni^2+} < \ce{Cu^2+} > \ce{Zn^2+}$. Seng ($d^{10}$) tidak memiliki stabilisasi medan ligan dan tidak mengalami
distorsi Jahn--Teller: seng kembali turun di bawah tembaga.
\end{solution}

\begin{solution}{exo:b3:bioinorganic:3}
$\ce{Ca^2+}$ dan $\ce{Fe^3+}$, keras: karboksilat (dan fenolat untuk besi). $\ce{Cu+}$ dan $\ce{Hg^2+}$, lunak:
tiolat sisteina. $\ce{Zn^2+}$, perbatasan: histidina dan sisteina.
\end{solution}

\begin{solution}{exo:b3:bioinorganic:4}
$\ce{2H2O -> O2 + 4H+ + 4e-}$: empat elektron, empat foton per $\ce{O2}$; empat mol foton per
mol $\ce{O2}$.
\end{solution}

\begin{solution}{exo:b3:bioinorganic:5}
Mioglobin: $\theta = 13/13.37 = 0.972$ dan $5/5.37 = 0.931$: mioglobin melepaskan $(0.972 - 0.931)/0.972 = 4$\,\% dari
oksigennya. Hemoglobin: $\theta = 0.972$ dan $0.724$: hemoglobin melepaskan 26\,\%.
\end{solution}

\begin{solution}{exo:b3:bioinorganic:6}
$[\mathrm{HbCO}]/[\mathrm{HbO_2}] = 230 \times \num{50e-6}/0.2095 = 0.055$; fraksinya $0.055/1.055 = 5.2$\,\%.
\end{solution}

\begin{solution}{exo:b3:bioinorganic:7}
Sulfida masing-masing $-2$, sisteinat masing-masing $-1$. $\ce{[Fe2S2(SR)4]^2-}$: jumlah bilangan oksidasi besi $-2 + 4 + 4 = +6$: dua
Fe(III). $\ce{[Fe2S2(SR)4]^3-}$: $+5$, satu Fe(III) dan satu Fe(II) (valensi campuran). $\ce{[Fe4S4(SR)4]^2-}$:
$-2 + 8 + 4 = +10$, dua Fe(III) dan dua Fe(II), dengan rata-rata $+2.5$ untuk setiap besi.
\end{solution}

\begin{solution}{exo:b3:bioinorganic:8}
$\ce{[PtCl2(NH3)2] + H2O -> [PtCl(H2O)(NH3)2]+ + Cl-}$. Konsentrasi klorida darah yang tinggi
mendorong kesetimbangan kembali; di dalam sel, kadar klorida jauh lebih rendah, dan kompleks
akua terbentuk. Platina mengikat N7 guanina. Pada isomer \textit{trans}, kedua
situs \omterm{def:b3:complex-mechanisms:labile}{labil} saling berseberangan dan tidak dapat mencapai dua basa bersebelahan pada satu
untai.
\end{solution}

\begin{solution}{exo:b3:bioinorganic:9}
$1/T_1 = 1/\qty{1.2}{s} + \qty{4.0}{L.mmol^{-1}.s^{-1}} \times \qty{0.10}{mmol/L} = 0.833 + 0.400 = \qty{1.23}{s^{-1}}$: $T_1 = \qty{0.81}{s}$.
\end{solution}

\begin{solution}{exo:b3:bioinorganic:10}
$K = K_{\mathrm{PbL}}/K_{\mathrm{CaL}} = 10^{18.0 - 10.7} = 10^{7.3} = \num{2e7}$: timbal menggeser kalsium sepenuhnya. Ligan
bebas juga akan mengikat kalsium darah dan menurunkan konsentrasinya secara
berbahaya; kompleks kalsium bertukar dengan timbal dan membiarkan kadar kalsium
tidak berubah.
\end{solution}

\begin{solution}{exo:b3:bioinorganic:11}
Jauh di bawah $K_M$, reaksinya berorde satu: $k' = (k_{\mathrm{cat}}/K_M)[\mathrm E]$, $t_{1/2} = \ln 2/k'$. Karbonat
anhidrase: $k' = \qty{100}{s^{-1}}$, $t_{1/2} = \qty{6.9}{ms}$. \omterm{def:b3:complex-kinetics:enzyme}{Enzim} yang khas: $k' = \qty{0.10}{s^{-1}}$, $t_{1/2} = \qty{6.9}{s}$, seribu
kali lebih lama, terlalu lambat untuk satu detik yang dilewatkan sel darah merah di kapiler paru-paru.
\end{solution}

\begin{solution}{exo:b3:bioinorganic:12}
Setengah situs ditempati CO berarti $[\mathrm{HbCO}] = [\mathrm{HbO_2}]$: $p(\ce{CO}) = p(\ce{O2})/M = 0.2095/230 = \qty{9.1e-4}{bar}$, sekitar
\qty{910}{ppm}.
\end{solution}

\begin{solution}{pb:b3:bioinorganic:1}
\textbf{1.} $\theta = 13^{2.7}/(3.5^{2.7} + 13^{2.7}) = 0.972$.

\textbf{2.} $\theta(\qty{5.0}{kPa}) = 0.724$.

\textbf{3.} $5.0/(0.37 + 5.0) = 0.931$.

\textbf{4.} Nilai $p_{50}$ mioglobin sepuluh kali lebih rendah: pada tekanan di otot yang sedang bekerja, mioglobin jauh lebih
jenuh daripada hemoglobin pada tekanan yang sama, sehingga oksigen berpindah dari
hemoglobin ke mioglobin.

\textbf{5.} $n = 1$.

\textbf{6.} Pengikatan pada satu \omterm{def:b3:bioinorganic:haem}{hem} menggerakkan besinya ke dalam bidang dan menarik histidina beserta
heliksnya; subunit-subunit bergeser bersama ke bentuk yang afinitasnya lebih tinggi untuk
situs-situs lainnya.

\textbf{7.} $\qty{150}{g/L}/\qty{64500}{g/mol} = \qty{2.33e-3}{mol/L}$.

\textbf{8.} $4 \times \num{2.33e-3} = \qty{9.30e-3}{mol/L}$.

\textbf{9.} $\qty{9.30}{mmol/L} \times 0.972 = \qty{9.04}{mmol}$.

\textbf{10.} $\qty{9.30}{mmol/L} \times (0.972 - 0.724) = \qty{2.31}{mmol}$.

\textbf{11.} $\qty{2.31e-3}{mol} \times \qty{22.41}{L/mol} = \qty{52}{mL}$.

\textbf{12.} $0.248/0.972 = 26$\,\%.

\textbf{13.} $5.0^{2.7}/(4.0^{2.7} + 5.0^{2.7}) = 0.646$.

\textbf{14.} $\qty{9.30}{mmol/L} \times (0.972 - 0.646) = \qty{3.03}{mmol}$.

\textbf{15.} $3.03/2.31 = 1.31$: tiga puluh persen lebih banyak.

\textbf{16.} Karbon dioksida hasil respirasi, yang dihidrasi menjadi asam karbonat oleh
karbonat anhidrase, dan asam laktat pada kerja berat.

\textbf{17.} $\qty{3.03}{mmol/L} \times \qty{5.0}{L/min} = \qty{15.1}{mmol/min}$.

\textbf{18.} $\qty{15.1e-3}{mol/min} \times \qty{22.41}{L/mol} = \qty{0.34}{L/min}$.

\textbf{19.} $230 \times \num{100e-6}/0.2095 = 0.110$.

\textbf{20.} $0.110/1.110 = 0.099$: sekitar 10\,\% situs.

\textbf{21.} Situs-situs yang masih membawa oksigen, dalam molekul yang juga membawa CO, berada dalam
bentuk berafinitas tinggi: kurvanya bergeser ke kiri dan situs-situs itu melepaskan lebih sedikit
oksigennya di jaringan.

\textbf{22.} Menurut hubungan persaingan, rasio $[\mathrm{HbCO}]/[\mathrm{HbO_2}]$ turun sebanding dengan $p(\ce{O2})$:
menghirup oksigen murni, yang tekanannya hampir lima kali tekanannya di udara, menggeser karbon
monoksida lebih cepat.

\textbf{23.} Dengan menghalangi pengikatan CO yang linear dan membentuk ikatan hidrogen dengan $\ce{O2}$, histidina itu menjaga $M$ tetap
bernilai ratusan alih-alih nilai yang jauh lebih besar pada \omterm{def:b3:bioinorganic:haem}{hem} bebas; jika tidak, bahkan
karbon monoksida yang dibuat tubuh sendiri akan menyumbat sebagian besar hemoglobin.

\textbf{24.} Saat istirahat, darah mengantarkan sekitar \qty{0.34}{L} oksigen per menit.
\end{solution}
